\section{Introduction}\label{sec:intro}

The thesis studies how paraglider pilots move in cross-country flight: how long they climb, how far
they glide, how their displacement grows with time. Flying in a group could change all of this.
In soaring birds, social information on the position of thermals changes the gliding strategy, and
in a migrating flock a few individuals find the thermals while the others follow them. The note
\emph{Group flights in free flight} (\texttt{ideas/group-identification}) labels every instant of
every flight of the FFVL archive as solo or in a group. This report asks whether, on the same
archive, pilots move differently when they fly in a group.

A raw comparison cannot answer, because group time differs from solo flight in four other ways.
\begin{enumerate}[label=(\roman*)]
  \item \textbf{The pilots.} Those who fly in groups may be more experienced, or fly faster wings.
  \item \textbf{The place.} Groups form on some sites and routes, where any pilot may progress
        straighter.
  \item \textbf{The day and the hour.} Group time gathers around midday and on some days more than
        others, and the air of that day and hour acts on every pilot.
  \item \textbf{The moment.} Pilots may join others when the progress is straight anyway, along a
        convergence line or a ridge, and the straight progress would then come before the group.
\end{enumerate}
Each of these produces a difference between group and solo time with no effect of the group. The
report removes them one at a time. Section~\ref{sec:setup} fixes the populations and the measures,
and Section~\ref{sec:raw} shows the raw difference. Contrasts within strata of site, day and hour,
and within the same flight, separate the pilots and the day from the group
(Section~\ref{sec:other}); reweighting the solo flights to the zones and hours where groups fly
tests the place (Section~\ref{sec:zones}); the same flight before, during and after a group tests
the moment (Section~\ref{sec:prepost}). What survives the three tests we call the group effect.
Section~\ref{sec:inside} looks inside the group, at its size and at who leads, and
Section~\ref{sec:conclusions} says what can and cannot be concluded.
